\paragraph{Warm start at the true values (Table~\ref{tab:warm}).} Initialising $Q$ at $Q^\ast$ removes the transient from zero initialisation. On the trap MDP the final $Q$-bias of Q-learning is \rhval{abl_warm/q-learning/maxbias/qbias_final/mean} (Q-learning has pulled $Q(A,\text{left})$ well above its true value of $-0.1$ even though it started at the truth), whereas Double's is \rhval{abl_warm/double-q-learning/maxbias/qbias_final/mean}, and the suboptimal fraction falls to \rhval{abl_warm/double-q-learning/maxbias/subopt_all/mean} versus \rhval{abl_warm/q-learning/maxbias/subopt_all/mean} ($p=$\rhval{cmp/abl_warm/q-learning/maxbias/subopt_all/welch_p}). This is the clean version of H1: Double nearly removes the overestimation, Q-learning does not. On random20 with warm start, the first-20-episode start-state bias is slightly positive for Q-learning (\rhval{abl_warm/q-learning/random20/bias_first/mean}), but by the end both estimators are below the truth (Q-learning \rhval{abl_warm/q-learning/random20/bias_final/mean}, Double \rhval{abl_warm/double-q-learning/random20/bias_final/mean}), so the negative final bias on random MDPs is not only a transient from the initial values. Notably the action ranking flips relative to the zero-initialised main run: with warm start Double takes fewer suboptimal actions than Q-learning (\rhval{abl_warm/double-q-learning/random20/subopt_all/mean} versus \rhval{abl_warm/q-learning/random20/subopt_all/mean}, $p=$\rhval{cmp/abl_warm/q-learning/random20/subopt_all/welch_p}). The cost of Double in the main random20 run is thus a cost of reaching the true values from a bad initialisation, not a property of the steady state.

\begin{table*}[t]\centering
\caption{Warm-start ablation (initial $Q=Q^\ast$), mean $\pm$ std over 5 seeds. Same metrics and conventions as Table~\ref{tab:main}.}\label{tab:warm}
\resizebox{0.9\textwidth}{!}{\begin{tabular}{llccccc}
\toprule
Method & Task & bias\_first $\downarrow$ & bias\_final $\downarrow$ & qbias\_final $\downarrow$ & subopt\_all $\downarrow$ & $n$ \\
\midrule
Q-learning & random20 & 0.003 $\pm$ 0.003 & -0.053 $\pm$ 0.005 & -0.098 $\pm$ 0.001 & 0.164 $\pm$ 0.002 & 5 \\
\textbf{Double Q-learning} & random20 & \underline{0.001 $\pm$ 0.002} & \underline{-0.100 $\pm$ 0.006} & \underline{-0.107 $\pm$ 0.002} & \underline{0.144 $\pm$ 0.002} & 5 \\
Maxmin Q-learning & random20 & \textbf{-0.120 $\pm$ 0.002} & \textbf{-0.507 $\pm$ 0.002} & \textbf{-0.464 $\pm$ 0.001} & 0.164 $\pm$ 0.002 & 5 \\
Weighted Double Q-learning & random20 & 0.002 $\pm$ 0.001 & -0.069 $\pm$ 0.006 & -0.078 $\pm$ 0.002 & \textbf{0.141 $\pm$ 0.001} & 5 \\
\midrule
Q-learning & maxbias & \textbf{0.000 $\pm$ 0.000} & 0.011 $\pm$ 0.001 & 0.075 $\pm$ 0.001 & 0.101 $\pm$ 0.002 & 5 \\
\textbf{Double Q-learning} & maxbias & \underline{0.000 $\pm$ 0.000} & 0.000 $\pm$ 0.000 & \underline{0.007 $\pm$ 0.001} & 0.052 $\pm$ 0.000 & 5 \\
Maxmin Q-learning & maxbias & 0.000 $\pm$ 0.000 & \textbf{0.000 $\pm$ 0.000} & \textbf{0.004 $\pm$ 0.000} & \textbf{0.050 $\pm$ 0.000} & 5 \\
Weighted Double Q-learning & maxbias & 0.000 $\pm$ 0.000 & \underline{0.000 $\pm$ 0.000} & 0.009 $\pm$ 0.001 & \underline{0.052 $\pm$ 0.000} & 5 \\
\bottomrule
\end{tabular}
}
\end{table*}

\paragraph{Simultaneous double update (Table~\ref{tab:both}).} Updating both tables on every sample with the cross targets gives, on the trap MDP, exactly the Q-learning numbers (peak bias \rhval{abl_double/double-q-learning-both-tables/maxbias/bias_peak/mean} versus \rhval{main/q-learning/maxbias/bias_peak/mean}; suboptimal fraction \rhval{abl_double/double-q-learning-both-tables/maxbias/subopt_all/mean} versus \rhval{main/q-learning/maxbias/subopt_all/mean}), and on random20 nearly the same (final bias \rhval{abl_double/double-q-learning-both-tables/random20/bias_final/mean} versus \rhval{main/q-learning/random20/bias_final/mean}). This follows from the symmetry argument of Section~\ref{sec:method}: cross-evaluation only helps when the two tables are fitted on disjoint samples. The small random20 differences come from the in-place update order, which we did not verify beyond this reasoning.

\begin{table*}[t]\centering
\caption{Simultaneous double update (both tables updated on each sample), mean $\pm$ std over 5 seeds.}\label{tab:both}
\resizebox{0.7\textwidth}{!}{\begin{tabular}{llcccc}
\toprule
Method & Task & bias\_peak $\downarrow$ & bias\_final $\downarrow$ & subopt\_all $\downarrow$ & $n$ \\
\midrule
Double Q-learning (both tables) & maxbias & \textbf{0.106 $\pm$ 0.002} & \textbf{0.008 $\pm$ 0.002} & \textbf{0.380 $\pm$ 0.005} & 5 \\
\midrule
Double Q-learning (both tables) & random20 & \textbf{-0.133 $\pm$ 0.011} & \textbf{-0.147 $\pm$ 0.012} & \textbf{0.255 $\pm$ 0.002} & 5 \\
\bottomrule
\end{tabular}
}
\end{table*}

\paragraph{Step size, noise and exploration on random20 (Table~\ref{tab:sweeps}, Fig.~\ref{fig:sweeps}).} The step-size sweep suggests that much of Double's penalty on random20 is a per-table learning-speed effect: Q-learning with half the step size ($\alpha=0.05$) has final bias \rhval{analysis/sweep-aggregation/maxbias/alpha_0p05_q_bias_final/mean} and suboptimal fraction \rhval{analysis/sweep-aggregation/maxbias/alpha_0p05_q_subopt_all/mean}, close to Double at $\alpha=0.1$ (\rhval{analysis/sweep-aggregation/maxbias/alpha_0p1_double_bias_final/mean} and \rhval{analysis/sweep-aggregation/maxbias/alpha_0p1_double_subopt_all/mean}); at $\alpha=0.2$ the gap in action quality shrinks (Q-learning \rhval{analysis/sweep-aggregation/maxbias/alpha_0p2_q_subopt_all/mean}, Double \rhval{analysis/sweep-aggregation/maxbias/alpha_0p2_double_subopt_all/mean}). This is consistent with, but does not prove, an effective-step-size explanation (each Double table sees half the updates). The noise sweep does \emph{not} support the first part of H4: Q-learning's final bias is not increasing in $\sigma$ at the grid points (\rhval{analysis/sweep-aggregation/maxbias/sigma_0p25_q_bias_final/mean}, \rhval{analysis/sweep-aggregation/maxbias/sigma_1p0_q_bias_final/mean}, \rhval{analysis/sweep-aggregation/maxbias/sigma_2p0_q_bias_final/mean}, \rhval{analysis/sweep-aggregation/maxbias/sigma_4p0_q_bias_final/mean} for $\sigma=0.25,1,2,4$), and more noise makes both estimators more negative at $\sigma=4$. The ranking of action quality does change with noise: Q-learning is better at $\sigma\le 2$ (suboptimal fraction \rhval{analysis/sweep-aggregation/maxbias/sigma_2p0_q_subopt_all/mean} versus \rhval{analysis/sweep-aggregation/maxbias/sigma_2p0_double_subopt_all/mean} at $\sigma=2$) but Double is better at $\sigma=4$ (\rhval{analysis/sweep-aggregation/maxbias/sigma_4p0_double_subopt_all/mean} versus \rhval{analysis/sweep-aggregation/maxbias/sigma_4p0_q_subopt_all/mean}), the only swept regime of random20 where the double estimator takes fewer suboptimal actions.

\paragraph{Exploration and adaptive sampling.} With $\epsilon=1$ (uniformly random behaviour, so the data do not depend on the estimates) Q-learning's final bias is positive (\rhval{analysis/sweep-aggregation/maxbias/epsilon_1p0_q_bias_final/mean}) while Double's is negative (\rhval{analysis/sweep-aggregation/maxbias/epsilon_1p0_double_bias_final/mean}): this is the textbook sign pattern. With $\epsilon=0.1$ both are negative, and the negativity shrinks as $\epsilon$ grows (Q-learning: \rhval{analysis/sweep-aggregation/maxbias/epsilon_0p1_q_bias_final/mean}, \rhval{analysis/sweep-aggregation/maxbias/epsilon_0p3_q_bias_final/mean}, \rhval{analysis/sweep-aggregation/maxbias/epsilon_1p0_q_bias_final/mean} for $\epsilon=0.1,0.3,1$). This pattern is consistent with the negative bias on random MDPs being caused by greedy, adaptively collected data (an action that looks bad early stops being sampled), but $\epsilon$ is the only factor we varied here, so we state it as a possibility rather than a finding.

\begin{table*}[t]\centering
\caption{Sensitivity on random20 (means over 5 seeds). Q = Q-learning, DQ = Double Q-learning; $\hat V$-bias: final start-state estimated-minus-true; $\hat Q$-bias: final all-pair $Q$-bias; subopt: suboptimal fraction over all episodes. The default setting of each sweep ($\sigma=1$, $\epsilon=0.1$, $\alpha=0.1$) is the main-run value.}\label{tab:sweeps}
\resizebox{0.9\textwidth}{!}{\begin{tabular}{llrrrrrr}\hline
Sweep & value & $\hat V$-bias Q & $\hat V$-bias DQ & $\hat Q$-bias Q & $\hat Q$-bias DQ & subopt Q & subopt DQ \\\hline
sigma & 0.25 & -0.181 & -0.433 & -0.452 & -0.876 & 0.276 & 0.348 \\
 & 1.0 & -0.150 & -0.364 & -0.444 & -0.840 & 0.255 & 0.312 \\
 & 2.0 & -0.209 & -0.452 & -0.563 & -0.924 & 0.278 & 0.298 \\
 & 4.0 & -0.527 & -1.156 & -1.211 & -1.658 & 0.377 & 0.335 \\
\hline
epsilon & 0.1 & -0.150 & -0.364 & -0.444 & -0.840 & 0.255 & 0.312 \\
 & 0.3 & -0.037 & -0.126 & -0.087 & -0.255 & 0.329 & 0.344 \\
 & 1.0 & 0.090 & -0.093 & 0.068 & -0.114 & 0.750 & 0.750 \\
\hline
alpha & 0.02 & -0.699 & -0.966 & -1.377 & -1.680 & 0.378 & 0.388 \\
 & 0.05 & -0.342 & -0.622 & -0.829 & -1.255 & 0.317 & 0.368 \\
 & 0.1 & -0.150 & -0.364 & -0.444 & -0.840 & 0.255 & 0.312 \\
 & 0.2 & -0.123 & -0.288 & -0.307 & -0.560 & 0.235 & 0.252 \\
\hline\end{tabular}
}
\end{table*}

\begin{figure*}[t]\centering\includegraphics[width=\textwidth]{sweeps.pdf}
\caption{Sensitivity sweeps (means over 5 seeds). First three panels: random20 final start-state bias of Q-learning (orange) and Double (green) versus $\sigma$, $\epsilon$, $\alpha$. Right: trap MDP peak start-state overestimate versus number of actions $n$ at state B for all four systems (orange Q, green Double, purple Weighted Double, grey Maxmin).}\label{fig:sweeps}
\end{figure*}

\paragraph{Number of actions at the trap state (Table~\ref{tab:actions}).} Q-learning's peak overestimate grows with the number of actions, from \rhval{analysis/sweep-aggregation/maxbias/actions_2_q_bias_peak/mean} at $n=2$ to \rhval{analysis/sweep-aggregation/maxbias/actions_100_q_bias_peak/mean} at $n=100$ (all five grid values are in the table and the sequence increases), while Double stays between \rhval{analysis/sweep-aggregation/maxbias/actions_2_double_bias_peak/mean} and \rhval{analysis/sweep-aggregation/maxbias/actions_100_double_bias_peak/mean}; the second part of H4 is supported on the trap MDP at these grid points. In action quality the trap grows costly for Q-learning (suboptimal fraction \rhval{analysis/sweep-aggregation/maxbias/actions_100_q_subopt_all/mean} at $n=100$) and for Maxmin ($N=2$: \rhval{analysis/sweep-aggregation/maxbias/actions_100_maxmin_subopt_all/mean}, versus \rhval{analysis/sweep-aggregation/maxbias/actions_100_double_subopt_all/mean} for Double). Maxmin's bias is small but its action quality degrades with $n$, a failure case of the two-estimator minimum that we did not investigate further.

\begin{table}[t]\centering
\caption{Trap MDP with $n$ actions at state B (means over 5 seeds). Q, DQ, WDQ, MM: Q-learning, Double, Weighted Double, Maxmin. Left block: peak start-state overestimate; right block: suboptimal fraction over all episodes.}\label{tab:actions}
\resizebox{\columnwidth}{!}{\begin{tabular}{rrrrrrrrr}\hline
 & \multicolumn{4}{c}{start-state peak bias} & \multicolumn{4}{c}{suboptimal fraction} \\
$n$ & Q & DQ & WDQ & MM & Q & DQ & WDQ & MM \\\hline
2 & 0.037 & 0.006 & 0.006 & 0.002 & 0.142 & 0.103 & 0.106 & 0.082 \\
5 & 0.077 & 0.008 & 0.008 & 0.007 & 0.246 & 0.111 & 0.117 & 0.126 \\
10 & 0.106 & 0.008 & 0.009 & 0.011 & 0.380 & 0.119 & 0.127 & 0.197 \\
25 & 0.123 & 0.009 & 0.009 & 0.021 & 0.675 & 0.122 & 0.133 & 0.380 \\
100 & 0.128 & 0.010 & 0.010 & 0.031 & 0.940 & 0.128 & 0.137 & 0.794 \\
\hline\end{tabular}
}
\end{table}
